\section{The formula}\label{Chap1}

\subsection{The Kontsevich matrix model with arbitrary potential}

Let $\mathcal{H}_N$ be the space of Hermitian $N \times N$ matrices
equipped with the Lebesgue measure
\[
\dd H= \prod_{i=1}^N \dd H_{ii}
 \prod_{1 \leq i < j \leq N} \dd\operatorname{Re}(H_{ij})\, \dd\operatorname{Im}(H_{ij}).
\]
Given a positive matrix $\Lambda = \operatorname{diag}(\lambda_1,\dots,\lambda_N)$,
we introduce the Gaussian probability measure on $\mathcal{H}_N$
\begin{equation}
\label{theGauss}
\dd \mathbb{P}_N(H) = \frac{\sqrt{\Delta(\boldsymbol{\lambda},\boldsymbol{\lambda})}}{2^{\frac{N}{2}} (2\pi)^{\frac{N^2}{2}}}\,\dd H {\rm e}^{-\frac{1}{2}{\rm Tr}(\Lambda H^2)}.
\end{equation}
We use the notations $\Delta(\boldsymbol{\lambda}) = \prod_{1 \leq i < j \leq N} (\lambda_j - \lambda_i)$ and $\Delta(\boldsymbol{\lambda},\boldsymbol{\mu}) = \prod_{1 \leq i,j \leq N} (\lambda_i + \mu_j)$. We denote $[n] = \{1,\dots,n\}$ and $\lambda_{\min} = \min\{\lambda_i\mid i \in [N]\}$.

Let $V_0$ be a continuous function on $\mathbb{R}$ such that the measure ${\rm e}^{- \frac{1}{2}\lambda_{\min} x^2 + V_0(x)}\,\dd x$ has finite moments on $\mathbb{R}$ (take for instance $V_0$ to be a polynomial of even degree with negative top coefficient). Then the measure $\dd \mathbb{P}_N(H) {\rm e}^{ \operatorname{Tr}V_0(H)}$ on $\mathcal{H}_N$ is finite. Let
\[
V_{\mathbf{t}}(x) = V_0(x) + \sum_{k \geq 0} t_{2k + 1} x^{2k + 1},
\]
where $\mathbf{t} = (t_{2k + 1})_{k \geq 1}$ are formal parameters. The partition function of the Kontsevich model with arbitrary potential is defined by
\begin{equation}
\label{Zdef} Z_N(\mathbf{t}) = \int_{\mathcal{H}_N} \dd\mathbb{P}_N(H) {\rm e}^{\operatorname{Tr}V_{\mathbf{t}}(H)}.
\end{equation}
This note aims at showing that $Z_N(\mathbf{t})$ is a
tau-function of the BKP hierarchy \cite{Date:1981qz}.

\subsection{Pfaffian formula}

First, we establish a Pfaffian formula for $Z_N(\mathbf{t})$. The proof proposed in Section~\ref{S2} consists in classical algebraic manipulations with matrix integrals and an analysis argument -- that one may find of independent interest, see Lemma~\ref{HUHH} and Remark~\ref{therem} -- to justify the extension of de~Bruijn's Pfaffian formula \cite{deBruijn1955} to singular kernels like the one appearing in \eqref{BZ}.
\begin{Theorem}
\label{th:1} For even $N$, we have
\begin{equation}
\label{BZ} Z_N(\mathbf{t}) = \frac{\sqrt{\Delta(\boldsymbol{\lambda},\boldsymbol{\lambda})}}{2^{\frac{N^2}{2}} (2\pi)^{\frac{N}{2}}\prod_{n = 1}^{N - 1} n!} \operatorname{Pf}_{0 \leq m,n \leq N - 1}(K_{N;m,n}(\mathbf{t})),
\end{equation}
where $\fint = \lim_{\epsilon \rightarrow 0} \int_{|x + y| \geq \epsilon}$ is the Cauchy principal value integral and
\begin{gather*}
 K_{N;m,n}(\mathbf{t}) = \fint_{\mathbb{R}^2} \frac{x - y}{x + y} F_{N;m}(x)F_{N;n}(y) {\rm e}^{V_{\mathbf{t}}(x) + V_{\mathbf{t}}(y)}\,\dd x \dd y, \\
F_{N;n}(x) = x^{2n} + \frac{(-2)^n n!}{\Delta(\boldsymbol{\lambda})}\det\bigl(\lambda_i^0 \big| \lambda_i^1 \big| \cdots \big| \lambda_i^{n - 1} \big| R_N\bigl(-\tfrac{1}{2}\lambda_ix^2\bigr) \big| \lambda_i^{n + 1} \big| \cdots \big| \lambda_i^{N - 1}\bigr), \\
R_N(\xi) = \frac{\xi^N}{(N - 1)!} \int_{0}^{1} \dd u\, (1 - u)^{N - 1} {\rm e}^{\xi u} = {\rm e}^{\xi}\biggl(1 - \frac{\Gamma(N;\xi)}{(N - 1)!}\biggr).
\end{gather*}
\end{Theorem}

Note that $R_N(\xi)$ is simply the $N$-th remainder of the Taylor series of ${\rm e}^{\xi}$. The expression given for $F_{N;n}(x)$ emphasises that it is $x^{2n} + O\bigl(x^{2N + 2}\bigr)$ as $x \rightarrow 0$, but we also have the equivalent expression
\[
F_{N;n}(x) = \frac{(-2)^n n!}{\Delta(\boldsymbol{\lambda})} \det\bigl(\lambda_i^0 \big| \lambda_i^1 \big| \cdots \big| \lambda_i^{n - 1} \big| {\rm e}^{-\frac{1}{2}\lambda_i x^2} \big| \lambda_i^{n + 1} \big| \cdots \big| \lambda_i^{N - 1}\bigr).
\]
For instance, the formula for $N = 2$ involves the two functions
\[
F_{2;0}(x) = \frac{{\rm e}^{-\frac{1}{2}\lambda_1 x^2} \lambda_2 - {\rm e}^{-\frac{1}{2}\lambda_2 x^2}\lambda_1}{\lambda_2 - \lambda_1}, \qquad F_{2;1}(x) = 2 \frac{{\rm e}^{-\frac{1}{2}\lambda_1 x^2} - {\rm e}^{-\frac{1}{2}\lambda_2 x^2}}{\lambda_2 - \lambda_1} .
\]

